%! Dimensions of the Four Subspaces — Unit 3, Lesson 3.5 — Group Activity (Answer Key)
\documentclass[10pt]{article}
\usepackage{linalg-article}
\usepackage{linalg-key}   % pulls in -boxes; do NOT also load -boxes
\newcommand{\TallMath}[1]{$\displaystyle #1\rule[-1.4em]{0pt}{3.2em}$}
\newcommand{\vv}{\mathbf{v}}
\newcommand{\ww}{\mathbf{w}}
\newcommand{\xx}{\mathbf{x}}
\newcommand{\yy}{\mathbf{y}}
\newcommand{\svec}{\mathbf{s}}
\newcommand{\zero}{\mathbf{0}}
\newcommand{\T}{\mathsf{T}}

\begin{document}

\pageheader{Unit 3, Lesson 3.5}{Group Activity --- Answer Key}
\namepartnerperiod

\vspace{0.10in}

\begin{headlinebox}{blush}
\small\textbf{The Four Subspaces.} An $m\times n$ matrix $A$ of rank $r$ has four fundamental subspaces. In
the \emph{input room} $\mathbb{R}^n$: the row space $C(A^{\T})$ ($\dim r$) and the nullspace $N(A)$
($\dim n-r$). In the \emph{output room} $\mathbb{R}^m$: the column space $C(A)$ ($\dim r$) and the left
nullspace $N(A^{\T})$ ($\dim m-r$). \emph{One} reduction of $A$ gives a basis for all four. Work with your
partner and \emph{justify} every answer.
\end{headlinebox}

\vspace{0.10in}

\begin{tcolorbox}[colback=white, colframe=black!40, title=\textbf{Tier R --- Remediate},
    fonttitle=\bfseries, arc=1.5mm, left=3mm, right=3mm, top=2mm, bottom=2mm, breakable]
The matrix below is already reduced: $R=\begin{bmatrix} 1 & 0 & 2 \\ 0 & 1 & 4 \\ 0 & 0 & 0 \end{bmatrix}$.
\begin{enumerate}[leftmargin=1.7em, itemsep=8pt]
  \item \textbf{Three numbers.} Rows $m=\ans{$3$}$\quad columns $n=\ans{$3$}$\quad pivots $r=\ans{$2$}$
  \item \textbf{Which room?} Circle the room each subspace lives in.
    \begin{enumerate}[label=(\alph*), leftmargin=1.8em, itemsep=4pt]
      \item Column space $C(A)$: \quad \ans{$\mathbb{R}^m$} / $\mathbb{R}^n$ \qquad\qquad
            (c)\; Nullspace $N(A)$: \quad $\mathbb{R}^m$ / \ans{$\mathbb{R}^n$}
      \item Row space $C(A^{\T})$: \quad $\mathbb{R}^m$ / \ans{$\mathbb{R}^n$} \qquad\quad
            (d)\; Left nullspace $N(A^{\T})$: \quad \ans{$\mathbb{R}^m$} / $\mathbb{R}^n$
    \end{enumerate}
  \item \textbf{Four dimensions from $r$.}
        $\dim C(A)=\ans{$2$}$\quad $\dim C(A^{\T})=\ans{$2$}$\quad
        $\dim N(A)=\ans{$1$}$\quad $\dim N(A^{\T})=\ans{$1$}$
\end{enumerate}
\end{tcolorbox}

\begin{tcolorbox}[colback=white, colframe=black!40, title=\textbf{Tier A --- Approaching Proficiency},
    fonttitle=\bfseries, arc=1.5mm, left=3mm, right=3mm, top=2mm, bottom=2mm, breakable]
For each matrix: reduce, then give a basis \emph{and} dimension for all four subspaces, and check both
counting rules $r+(n-r)=n$ and $r+(m-r)=m$. Note that the two rooms $\mathbb{R}^m$ and $\mathbb{R}^n$ have
\emph{different} sizes here.
\begin{enumerate}[leftmargin=1.7em, itemsep=9pt]
  \item $A=\begin{bmatrix} 1 & 1 & 2 \\ 1 & 2 & 3 \end{bmatrix}$ \quad (\emph{$2\times3$; watch what the left nullspace does}).
        \ansline{Reduce to $\left[\begin{smallmatrix} 1&0&1\\0&1&1 \end{smallmatrix}\right]$: $m=2$, $n=3$, $r=2$. \textbf{$C(A)\subseteq\mathbb{R}^2$:} pivot columns $\begin{bsmallmatrix} 1\\1 \end{bsmallmatrix},\begin{bsmallmatrix} 1\\2 \end{bsmallmatrix}$, $\dim=2$ (all of $\mathbb{R}^2$). \textbf{$C(A^{\T})\subseteq\mathbb{R}^3$:} rows $(1,0,1),(0,1,1)$, $\dim=2$. \textbf{$N(A)\subseteq\mathbb{R}^3$:} $\svec=\begin{bsmallmatrix} -1\\-1\\1 \end{bsmallmatrix}$, $\dim=n-r=1$. \textbf{$N(A^{\T})\subseteq\mathbb{R}^2$:} $\dim=m-r=0$, so $N(A^{\T})=\{\zero\}$ (the rows are independent --- no row redundancy). Checks: $2+1=3=n$, $2+0=2=m$.}
  \item $A=\begin{bmatrix} 1 & 1 & 2 & 1 \\ 1 & 2 & 3 & 1 \\ 2 & 3 & 5 & 2 \end{bmatrix}$ \quad (\emph{$3\times4$; row~3 $=$ row~1 $+$ row~2}).
        \ansline{Reduce to $\left[\begin{smallmatrix} 1&0&1&1\\0&1&1&0\\0&0&0&0 \end{smallmatrix}\right]$: $m=3$, $n=4$, $r=2$. \textbf{$C(A)\subseteq\mathbb{R}^3$:} pivot columns $\begin{bsmallmatrix} 1\\1\\2 \end{bsmallmatrix},\begin{bsmallmatrix} 1\\2\\3 \end{bsmallmatrix}$, $\dim=2$. \textbf{$C(A^{\T})\subseteq\mathbb{R}^4$:} rows $(1,0,1,1),(0,1,1,0)$, $\dim=2$. \textbf{$N(A)\subseteq\mathbb{R}^4$:} $\svec_1=\begin{bsmallmatrix} -1\\-1\\1\\0 \end{bsmallmatrix},\svec_2=\begin{bsmallmatrix} -1\\0\\0\\1 \end{bsmallmatrix}$, $\dim=n-r=2$. \textbf{$N(A^{\T})\subseteq\mathbb{R}^3$:} row$_1+$row$_2-$row$_3=\zero$, so $\yy=\begin{bsmallmatrix} 1\\1\\-1 \end{bsmallmatrix}$, $\dim=m-r=1$. Checks: $2+2=4=n$, $2+1=3=m$.}
\end{enumerate}
\end{tcolorbox}

\begin{tcolorbox}[colback=white, colframe=black!40, title=\textbf{Tier E --- Extension},
    fonttitle=\bfseries, arc=1.5mm, left=3mm, right=3mm, top=2mm, bottom=2mm, breakable]
\textbf{Why the counting works.}
\begin{enumerate}[leftmargin=1.7em, itemsep=9pt]
  \item \textbf{Row rank $=$ column rank.} A reduction of $A$ has exactly $r$ pivots. Explain why that one
        number counts \emph{both} the independent columns (the pivot columns) and the independent rows (the
        nonzero rows of the reduced form) --- so $\dim C(A)=\dim C(A^{\T})=r$.
        \ansline{Each pivot sits in exactly one column (a pivot column) and exactly one row (a nonzero row of the reduced form). The $r$ pivot columns are an independent set spanning $C(A)$, and the $r$ nonzero reduced rows are an independent set spanning the row space. So the \emph{same} $r$ pivots count the independent columns and the independent rows: $\dim C(A)=\dim C(A^{\T})=r$.}
  \item \textbf{When is the left nullspace just $\{\zero\}$?} Explain why $N(A^{\T})=\{\zero\}$ exactly when
        $r=m$ (``full row rank''), and what that says about redundancy among the rows.
        \ansline{$\dim N(A^{\T})=m-r$, which is $0$ exactly when $r=m$. Then every row is a pivot row --- there are no zero rows and no combination of the rows gives $\zero$ except the trivial one, so the rows are independent (no row is redundant).}
  \item \textbf{Read the row dependency.} For $A=\begin{bmatrix} 1 & 1 & 2 \\ 2 & 1 & 3 \\ 3 & 2 & 5 \end{bmatrix}$
        a left-nullspace vector is $\yy=\begin{bsmallmatrix} 1\\1\\-1 \end{bsmallmatrix}$. Write the exact
        combination of the rows of $A$ that this says equals the zero row.
        \ansline{$\yy^{\T}A=\zero$ means $1\cdot\text{row}_1+1\cdot\text{row}_2-1\cdot\text{row}_3=\zero$, i.e.\ $\text{row}_3=\text{row}_1+\text{row}_2$. Check: $(1,1,2)+(2,1,3)=(3,2,5)$, the third row.}
\end{enumerate}
\end{tcolorbox}

\begin{teachernote}
Tier~A is the core skill: from one reduction, a basis and dimension for all four subspaces, with both
counting rules as the sanity check. Problem~1 is the ``full row rank'' surprise --- the left nullspace is
$\{\zero\}$ because $m-r=0$; use it to stress that a subspace \emph{can} be just the zero vector. Problem~2
uses $m\ne n$ so the two rooms ($\mathbb{R}^3$ outputs, $\mathbb{R}^4$ inputs) are visibly different sizes.
Tier~E~1 is the ``why row rank $=$ column rank'' argument (same pivots, counted two ways); if Tier~E runs
long, that is the must-do.
\end{teachernote}

\end{document}
